%%%%%%%%%%%%%%%%%%%%%%%%%%%%%%%%%%%%%%%%%%%%%%%%%%%%%%%%%%%%%%%%%%%%%%%%%%%%%%%%
%                                                                              %
% BlackArch Linux Guide (Persian / fa)                                         %
%                                                                              %
% Build with XeLaTeX (NOT pdflatex):  xelatex blackarch-guide-fa.tex           %
%                                                                              %
%%%%%%%%%%%%%%%%%%%%%%%%%%%%%%%%%%%%%%%%%%%%%%%%%%%%%%%%%%%%%%%%%%%%%%%%%%%%%%%%

\documentclass[a4paper, oneside, 11pt]{book}

%%% INCLUDES %%%
\usepackage{array}
\usepackage{xcolor}
\usepackage{enumerate}
\usepackage{fancyhdr}
\usepackage{fancyvrb}
\usepackage{geometry}
\usepackage{graphicx}
\usepackage{hyperref}
\usepackage{listings}
\usepackage{supertabular}
\usepackage{tocloft}
% xepersian must be loaded last (after hyperref)
\usepackage{xepersian}

%%% FONTS %%%
% Persian text font. Amiri ships with TeX Live (package "amiri") and covers all glyphs used here.
\settextfont{Amiri}
% Latin font (the original guide uses a sans-serif Latin font)
\setlatintextfont{Latin Modern Sans}

% the Persian font has no bullet glyph: take it from the Latin font
\renewcommand{\labelitemi}{\lr{\textbullet}}

%%% LAYOUT %%%
\setlength{\parindent}{0em}
\setlength{\parskip}{1.5ex plus0.5ex minus0.5ex}
\geometry{left=2.5cm, textwidth=16cm, top=3cm, textheight=25cm, bottom=3cm}
\widowpenalty=2000
\clubpenalty=1000
\frenchspacing
\sloppy
\pagecolor[HTML]{FFFFFF}
\color[HTML]{333333}
\setcounter{tocdepth}{10}
\setcounter{secnumdepth}{10}

\hypersetup{
  pdftitle={BlackArch Linux, The BlackArch Linux Guide},
  pdfsubject={BlackArch Linux, The BlackArch Linux Guide},
  pdfauthor={BlackArch Linux, BlackArch Linux},
  pdfkeywords={BlackArch Linux, Penetration Testing, Security, Hacking, Linux},
  pdfcenterwindow=true,
  colorlinks=true,
  breaklinks=true,
  linkcolor=blue,
  menucolor=blue,
  urlcolor=blue
}

% syntax highlighting
% all options should be set here document wide
\lstset{
backgroundcolor=\color[HTML]{EEEEEE},
frame=single,
basicstyle=\footnotesize\ttfamily,
float,
deletekeywords={return},
otherkeywords={mkdir, curl, sudo, sha1sum, grep, cut, sort, wget, makepkg,
pacman, blackman, chmod},
keywordstyle=\color{orange},
commentstyle=\color{blue},
stringstyle=\color{red},
language=bash,
showspaces=false,
showtabs=false,
tabsize=2
}

% NOTE: every lstlisting is wrapped in the xepersian `latin' environment;
% without it listings are not typeset in a right-to-left document.

%%% HEADER / FOOTER %%%
\setlength{\headheight}{33pt}
\setlength{\headsep}{33pt}
\lhead{{\includegraphics[width=1cm,height=1cm]{images/logo.png}}}
\rhead{راهنمای \lr{BlackArch Linux}}

%------------------%
%  TITLE PAGE      %
%------------------%
\begin{document}
\pagestyle{empty}
\begin{center}
\begin{figure}[htbp]
\centering
\vspace{0.5cm}
\includegraphics[width=8cm]{images/logo.png}
\label{fig:logo}
\end{figure}
\vspace{0.5cm}
\Huge{\textbf{راهنمای \lr{BlackArch Linux}}}\\
\vspace{1cm}
\Large{\color{blue}\lr{https://www.blackarch.org/}}\\
\vspace{0.5cm}
\end{center}
\newpage
\tableofcontents
\newpage
\pagestyle{fancy}

%------------------%
%  Chapter 1       %
%------------------%

\chapter{مقدمه}

\section{نگاه کلی}
راهنمای \lr{BlackArch Linux} از چند بخش تشکیل شده است:
\begin{itemize}
\item مقدمه: نگاهی کلی به پروژه، معرفی آن و اطلاعات مفید دیگر درباره‌ی آن
\item راهنمای کاربر: هر آنچه یک کاربر معمولی برای استفاده‌ی مؤثر از \lr{BlackArch} باید بداند
\item راهنمای توسعه‌دهنده: شروع توسعه و مشارکت در \lr{BlackArch}
\item راهنمای ابزارها: توضیح تفصیلی ابزارها همراه با نمونه‌های کاربرد (در دست تکمیل)
\end{itemize}

\section{\lr{BlackArch Linux} چیست؟}
\lr{BlackArch} یک توزیع کامل لینوکس برای متخصصان تست نفوذ و پژوهشگران امنیت است.
این توزیع از \href{https://www.archlinux.org/}{\lr{ArchLinux}} مشتق شده است و کاربران
می‌توانند اجزای \lr{BlackArch} را به‌صورت تکی یا گروهی، مستقیماً روی آن نصب کنند.

مجموعه‌ی ابزارها به‌شکل یک
\href{https://wiki.archlinux.org/index.php/Unofficial\_User\_Repositories}
{مخزن غیررسمی کاربران} برای \lr{Arch Linux} منتشر می‌شود؛ بنابراین می‌توانید
\lr{BlackArch} را روی یک نصب موجود \lr{Arch Linux} اضافه کنید. بسته‌ها را می‌توان تکی یا
بر اساس دسته نصب کرد.

این مخزن که پیوسته در حال گسترش است، هم‌اکنون بیش از
\href{https://www.blackarch.org/tools.html}{۲۶۰۰} ابزار دارد.
برای حفظ کیفیت مخزن، همه‌ی ابزارها پیش از افزوده‌شدن به پایگاه کد
به‌دقت آزمایش می‌شوند.
% should quickly describe the testing methods/code review procedures etc

\section{تاریخچه‌ی \lr{BlackArch Linux}}
به‌زودی...

\section{پلتفرم‌های پشتیبانی‌شده}
به‌زودی...

\section{به ما بپیوندید}
برای ارتباط با تیم \lr{BlackArch} می‌توانید از راه‌های زیر استفاده کنید:

وب‌سایت: \url{https://www.blackarch.org/}

ایمیل: \href{mailto:team@blackarch.org}{\lr{team@blackarch.org}}

\lr{IRC}: \url{irc://irc.freenode.net/blackarch}

\lr{Twitter}: \url{https://twitter.com/blackarchlinux}

\lr{GitHub}: \url{https://github.com/Blackarch/}

\lr{Discord}: \url{https://discord.com/invite/xMHt8dW}

%------------------%
%  Chapter 2       %
%------------------%


\chapter{راهنمای کاربر}

\section{نصب}
بخش‌های زیر نشان می‌دهند چگونه مخزن \lr{BlackArch} را راه‌اندازی کنید و بسته‌ها را
نصب کنید. \lr{BlackArch} از هر دو روش پشتیبانی می‌کند: نصب از مخزن با استفاده از
بسته‌های باینری، و کامپایل و نصب از روی سورس.

\lr{BlackArch} با نصب‌های معمولی \lr{Arch} سازگار است و مانند یک مخزن غیررسمی کاربران
عمل می‌کند. اگر به‌جای آن یک \lr{ISO} می‌خواهید، بخش
\href{https://www.blackarch.org/downloads.html#iso}{\lr{ISO}} را ببینید.

\subsection{نصب روی \lr{ArchLinux}}
اسکریپت \href{https://blackarch.org/strap.sh}{\lr{strap.sh}} را با دسترسی \lr{root} اجرا
کنید و دستورهای آن را دنبال کنید. نمونه‌ی زیر را ببینید.
\begin{latin}
\begin{lstlisting}
   curl -O https://blackarch.org/strap.sh
   sha1sum strap.sh # should match: 5ea40d49ecd14c2e024deecf90605426db97ea0c
   sudo chmod +x strap.sh
   sudo ./strap.sh
\end{lstlisting}
\end{latin}
% translator's addition: the English comment inside the listing is not translated
حاصل \lr{sha1sum} باید با مقداری که در کامنت همان خط آمده برابر باشد.

اکنون یک نسخه‌ی تازه از فهرست اصلی بسته‌ها را دریافت کنید و بسته‌ها را
همگام‌سازی کنید:
\begin{latin}
\begin{lstlisting}
  sudo pacman -Syyu
\end{lstlisting}
\end{latin}


\subsection{نصب بسته‌ها}
اکنون می‌توانید ابزارها را از مخزن \lr{blackarch} نصب کنید.
\begin{enumerate}
\item برای فهرست‌کردن همه‌ی ابزارهای موجود، این دستور را اجرا کنید:
\begin{latin}
\begin{lstlisting}
  pacman -Sgg | grep blackarch | cut -d' ' -f2 | sort -u
\end{lstlisting}
\end{latin}

\item برای نصب همه‌ی ابزارها، این دستور را اجرا کنید:
\begin{latin}
\begin{lstlisting}
  pacman -S blackarch
\end{lstlisting}
\end{latin}

\item برای نصب یک دسته از ابزارها، این دستور را اجرا کنید:
\begin{latin}
\begin{lstlisting}
  pacman -S blackarch-<category>
\end{lstlisting}
\end{latin}

\item برای دیدن دسته‌های \lr{blackarch}، این دستور را اجرا کنید:
\begin{latin}
\begin{lstlisting}
  pacman -Sg | grep blackarch
\end{lstlisting}
\end{latin}

\end{enumerate}

\subsection{نصب بسته‌ها از روی سورس}
به‌عنوان روشی جایگزین برای نصب، می‌توانید بسته‌های \lr{BlackArch} را از روی سورس
بسازید. فایل‌های \lr{PKGBUILD} را در
\href{https://github.com/BlackArch/blackarch/tree/master/packages}{\lr{GitHub}}
پیدا می‌کنید. برای ساخت کل مخزن می‌توانید از ابزار
\href{https://github.com/BlackArch/blackman}{\lr{Blackman}} استفاده کنید.
\begin{itemize}
\item ابتدا باید \lr{Blackman} را نصب کنید. اگر مخزن بسته‌های \lr{BlackArch} روی
دستگاه شما راه‌اندازی شده باشد، می‌توانید \lr{Blackman} را این‌گونه نصب کنید:
\begin{latin}
\begin{lstlisting}
  pacman -S blackman
\end{lstlisting}
\end{latin}

\item می‌توانید \lr{Blackman} را از روی سورس بسازید و نصب کنید:
\begin{latin}
\begin{lstlisting}
  mkdir blackman
  cd blackman
  wget https://raw.github.com/BlackArch/blackarch/master/packages/blackman/PKGBUILD
  # Make sure the PKGBUILD has not been maliciously tampered with.
  makepkg -s
\end{lstlisting}
\end{latin}
% translator's addition: the English comment inside the listing is not translated
پیش از اجرای \lr{makepkg} مطمئن شوید فایل \lr{PKGBUILD} دست‌کاری نشده است
(همان کامنت داخل دستورها).

\item یا \lr{Blackman} را از \lr{AUR} نصب کنید:
\begin{latin}
\begin{lstlisting}
  <whatever AUR helper you use> -S blackman
\end{lstlisting}
\end{latin}

\end{itemize}

\subsection{روش استفاده‌ی پایه از \lr{Blackman}} استفاده از \lr{Blackman} بسیار ساده است،
هرچند گزینه‌های آن با آنچه معمولاً از ابزاری مثل \lr{pacman} انتظار دارید فرق دارد.
روش‌های پایه‌ی استفاده در ادامه آمده است.
\begin{itemize}
\item دریافت، کامپایل و نصب بسته‌ها:
\begin{latin}
\begin{lstlisting}
  sudo blackman -i package
\end{lstlisting}
\end{latin}

\item دریافت، کامپایل و نصب یک دسته‌ی کامل:
\begin{latin}
\begin{lstlisting}
  sudo blackman -g group
\end{lstlisting}
\end{latin}

\item دریافت، کامپایل و نصب همه‌ی ابزارهای \lr{BlackArch}:
\begin{latin}
\begin{lstlisting}
  sudo blackman -a
\end{lstlisting}
\end{latin}

\item برای فهرست‌کردن دسته‌های \lr{blackarch}:
\begin{latin}
\begin{lstlisting}
  blackman -l
\end{lstlisting}
\end{latin}

\item برای فهرست‌کردن ابزارهای یک دسته:
\begin{latin}
\begin{lstlisting}
  blackman -p category
\end{lstlisting}
\end{latin}

\end{itemize}

\subsection{نصب از \lr{ISO} کامل، \lr{ISO} نصب شبکه‌ای (\lr{netinstall}) یا از \lr{ArchLinux}}
می‌توانید \lr{BlackArch Linux} را از یکی از \lr{ISO}های کامل (\lr{full}) یا \lr{netinstall} ما
نصب کنید.\\برای جزئیات به \url{https://www.blackarch.org/download.html#iso}
مراجعه کنید. پس از بوت‌شدن \lr{ISO}، مراحل زیر لازم است.
\begin{itemize}
\item بسته‌ی \lr{blackarch-installer} را نصب کنید:
\begin{latin}
\begin{lstlisting}
  sudo pacman -S blackarch-installer
\end{lstlisting}
\end{latin}

\item این دستور را اجرا کنید:
\begin{latin}
\begin{lstlisting}
  sudo blackarch-install
\end{lstlisting}
\end{latin}

\end{itemize}

%------------------%
%  Chapter 3       %
%------------------%

\chapter{راهنمای توسعه‌دهنده}

\section{سیستم ساخت و مخازن \lr{Arch}}

فایل‌های \lr{PKGBUILD} اسکریپت‌های ساخت هستند. هر کدام به \lr{makepkg(1)} می‌گوید
چگونه یک بسته بسازد. فایل‌های \lr{PKGBUILD} به زبان \lr{Bash} نوشته می‌شوند.

برای اطلاعات بیشتر، موارد زیر را بخوانید (یا دست‌کم مرورشان کنید):
\begin{itemize}
\item \href{https://wiki.archlinux.org/index.php/Creating_Packages}{\lr{Arch Wiki: Creating Packages}}
\item \href{https://wiki.archlinux.org/index.php/Makepkg}{\lr{Arch Wiki: makepkg}}
\item \href{https://wiki.archlinux.org/index.php/PKGBUILD}{\lr{Arch Wiki: PKGBUILD}}
\item \href{https://wiki.archlinux.org/index.php/Arch_Packaging_Standards}{\lr{Arch Wiki: Arch Packaging Standards}}
\end{itemize}

\section{استانداردهای \lr{PKGBUILD} در \lr{BlackArch}}
برای سادگی، فایل‌های \lr{PKGBUILD} ما شبیه نمونه‌های \lr{AUR} است و فقط چند تفاوت کوچک
دارد که در ادامه آمده است. هر بسته دست‌کم باید عضو گروه \lr{blackarch} باشد؛
هم‌پوشانی زیادی هم وجود خواهد داشت، یعنی بسیاری از بسته‌ها عضو چند گروه
می‌شوند.

\subsection{گروه‌ها}
برای اینکه کاربران بتوانند دسته‌ی مشخصی از بسته‌ها را سریع و آسان نصب کنند،
بسته‌ها در گروه‌هایی جدا شده‌اند. با گروه‌ها، کاربر فقط کافی است دستور
\lr{\texttt{pacman -S <group name>}} را اجرا کند تا تعداد زیادی بسته نصب شود.

\subsubsection{\lr{blackarch}}
گروه \lr{blackarch} گروه پایه‌ای است که همه‌ی بسته‌ها باید عضو آن باشند. به این ترتیب
کاربران می‌توانند به‌راحتی همه‌ی بسته‌ها را نصب کنند.

آنچه باید اینجا باشد: همه‌چیز.

\subsubsection{\lr{blackarch-anti-forensic}}
بسته‌هایی که برای مقابله با فعالیت‌های فارنزیک به کار می‌روند، از جمله
رمزنگاری، پنهان‌نگاری (\lr{steganography}) و هر چیزی که فایل‌ها یا ویژگی‌های
فایل‌ها را تغییر می‌دهد. این گروه به‌طور کلی شامل ابزارهایی است که برای
پنهان‌کردن اطلاعات، در سیستم تغییر ایجاد می‌کنند.

نمونه‌ها: \lr{luks, TrueCrypt, Timestomp, dd, ropeadope, secure-delete}

\subsubsection{\lr{blackarch-automation}}
بسته‌هایی که برای خودکارسازی ابزارها یا گردش‌کارها به کار می‌روند.

نمونه‌ها: \lr{blueranger, tiger, wiffy}

\subsubsection{\lr{blackarch-backdoor}}
بسته‌هایی که روی سیستم‌های از قبل آسیب‌پذیر، از درهای پشتی (\lr{backdoor})
بهره‌برداری می‌کنند یا درهای پشتی ایجاد می‌کنند.

نمونه‌ها: \lr{backdoor-factory, rrs, weevely}

\subsubsection{\lr{blackarch-binary}}
بسته‌هایی که به‌نوعی روی فایل‌های باینری کار می‌کنند.

نمونه‌ها: \lr{binwally, packerid}

\subsubsection{\lr{blackarch-bluetooth}}
بسته‌هایی که از هر چیزی مرتبط با استاندارد بلوتوث (\lr{802.15.1}) بهره‌برداری
می‌کنند.

نمونه‌ها: \lr{ubertooth, tbear, redfang}

\subsubsection{\lr{blackarch-code-audit}}
بسته‌هایی که برای تحلیل آسیب‌پذیری، کد منبع موجود را ممیزی می‌کنند.

نمونه‌ها: \lr{flawfinder, pscan}

\subsubsection{\lr{blackarch-cracker}}
بسته‌هایی که برای شکستن توابع رمزنگاری، یعنی هش‌ها، به کار می‌روند.

نمونه‌ها: \lr{hashcat, john, crunch}

\subsubsection{\lr{blackarch-crypto}}
بسته‌هایی که با رمزنگاری سروکار دارند، به‌جز شکستن رمز (\lr{cracking}).

نمونه‌ها: \lr{ciphertest, xortool, sbd}

\subsubsection{\lr{blackarch-database}}
بسته‌هایی که به هر سطحی از بهره‌برداری از پایگاه‌های داده مربوط می‌شوند.

نمونه‌ها: \lr{metacoretex, blindsql}

\subsubsection{\lr{blackarch-debugger}}
بسته‌هایی که به کاربر اجازه می‌دهند ببیند یک برنامه‌ی مشخص در همان لحظه
«چه می‌کند».

نمونه‌ها: \lr{radare2, shellnoob}

\subsubsection{\lr{blackarch-decompiler}}
بسته‌هایی که می‌کوشند برنامه‌ی کامپایل‌شده را به کد منبع برگردانند.

نمونه‌ها: \lr{flasm, jd-gui}

\subsubsection{\lr{blackarch-defensive}}
بسته‌هایی که برای محافظت از کاربر در برابر بدافزار و حمله‌ی دیگر کاربران به کار
می‌روند.

نمونه‌ها: \lr{arpon, chkrootkit, sniffjoke}

\subsubsection{\lr{blackarch-disassembler}}
این گروه شبیه \lr{blackarch-decompiler} است و احتمالاً برنامه‌های زیادی در هر دو
گروه قرار می‌گیرند؛ با این تفاوت که این بسته‌ها به‌جای کد منبع خام، خروجی
اسمبلی تولید می‌کنند.

نمونه‌ها: \lr{inguma, radare2}

\subsubsection{\lr{blackarch-dos}}
بسته‌هایی که از حمله‌های \lr{DoS} (انکار سرویس، \lr{Denial of Service}) استفاده می‌کنند.

نمونه‌ها: \lr{42zip, nkiller2}

\subsubsection{\lr{blackarch-drone}}
بسته‌هایی که برای مدیریت پهپادهای (\lr{drone}) فیزیکی به کار می‌روند.

نمونه‌ها: \lr{meshdeck, skyjack}

\subsubsection{\lr{blackarch-exploitation}}
بسته‌هایی که از اکسپلویت‌ها (\lr{exploit}) در برنامه‌ها یا سرویس‌های دیگر بهره
می‌گیرند.

نمونه‌ها: \lr{armitage, metasploit, zarp}

\subsubsection{\lr{blackarch-fingerprint}}
بسته‌هایی که از تجهیزات زیست‌سنجی اثر انگشت بهره‌برداری می‌کنند.

نمونه‌ها: \lr{dns-map, p0f, httprint}

\subsubsection{\lr{blackarch-firmware}}
بسته‌هایی که از آسیب‌پذیری‌های فریم‌ویر (\lr{firmware}) بهره‌برداری می‌کنند.

نمونه‌ها: هنوز موردی نیست؛ باید هرچه زودتر تکمیل شود.

\subsubsection{\lr{blackarch-forensic}}
بسته‌هایی که برای یافتن داده روی دیسک‌های فیزیکی یا حافظه‌های توکار
(\lr{embedded}) به کار می‌روند.

نمونه‌ها: \lr{aesfix, nfex, wyd}

\subsubsection{\lr{blackarch-fuzzer}}
بسته‌هایی که بر اصل آزمون فاز (\lr{fuzz testing}) کار می‌کنند، یعنی «پرتاب‌کردن»
ورودی‌های تصادفی به سمت هدف تا ببینیم چه اتفاقی می‌افتد.

نمونه‌ها: \lr{msf, mdk3, wfuzz}

\subsubsection{\lr{blackarch-hardware}}
بسته‌هایی که از هر چیزی مرتبط با سخت‌افزار فیزیکی بهره‌برداری می‌کنند یا آن را
مدیریت می‌کنند.

نمونه‌ها: \lr{arduino, smali}

\subsubsection{\lr{blackarch-honeypot}}
بسته‌هایی که نقش «هانی‌پات» (\lr{honeypot}) را دارند، یعنی برنامه‌هایی که شبیه
سرویس‌های آسیب‌پذیر به نظر می‌رسند و برای کشاندن هکرها به دام استفاده
می‌شوند.

نمونه‌ها: \lr{artillery, bluepot, wifi-honey}

\subsubsection{\lr{blackarch-keylogger}}
بسته‌هایی که ضربه‌های کلید را روی سیستمی دیگر ثبت و نگهداری می‌کنند.

نمونه‌ها: هنوز موردی نیست؛ باید هرچه زودتر تکمیل شود.

\subsubsection{\lr{blackarch-malware}}
بسته‌هایی که هر نوع نرم‌افزار مخرب به‌شمار می‌آیند یا برای تشخیص بدافزار به کار
می‌روند.

نمونه‌ها: \lr{malwaredetect, peepdf, yara}

\subsubsection{\lr{blackarch-misc}}
بسته‌هایی که به‌طور مشخص در هیچ‌یک از دسته‌ها نمی‌گنجند.

نمونه‌ها: \lr{oh-my-zsh-git, winexe, stompy}

\subsubsection{\lr{blackarch-mobile}}
بسته‌هایی که پلتفرم‌های موبایل را دست‌کاری می‌کنند.

نمونه‌ها: \lr{android-sdk-platform-tools, android-udev-rules}

\subsubsection{\lr{blackarch-networking}}
بسته‌هایی که به شبکه‌های \lr{IP} مربوط می‌شوند.

نمونه‌ها: \lr{arptools, dnsdiag, impacket}

\subsubsection{\lr{blackarch-nfc}}
بسته‌هایی که از \lr{NFC} (ارتباط میدان نزدیک، \lr{near-field communications}) استفاده
می‌کنند.

نمونه‌ها: \lr{nfcutils}

\subsubsection{\lr{blackarch-packer}}
بسته‌هایی که روی پکرها (\lr{packer}) کار می‌کنند یا به آنها مربوط می‌شوند.

\textit{پکرها برنامه‌هایی هستند که بدافزار را درون فایل‌های اجرایی دیگر جاسازی
می‌کنند.}

نمونه‌ها: \lr{packerid}

\subsubsection{\lr{blackarch-proxy}}
بسته‌هایی که نقش پراکسی (\lr{proxy}) دارند، یعنی ترافیک را از طریق گره‌ای دیگر در
اینترنت هدایت می‌کنند.

نمونه‌ها: \lr{burpsuite, ratproxy, sslnuke}

\subsubsection{\lr{blackarch-recon}}
بسته‌هایی که به‌طور فعال به دنبال آسیب‌پذیری‌های قابل بهره‌برداری در دنیای واقعی
(\lr{in the wild}) می‌گردند. این گروه بیشتر یک گروه چتری برای بسته‌های مشابه است.

نمونه‌ها: \lr{canri, dnsrecon, netmask}

\subsubsection{\lr{blackarch-reversing}}
این یک گروه چتری برای هر دیکامپایلر، دیس‌اسمبلر یا برنامه‌ی مشابه است.

نمونه‌ها: \lr{capstone, radare2, zerowine}

\subsubsection{\lr{blackarch-scanner}}
بسته‌هایی که سیستم‌های انتخاب‌شده را برای یافتن آسیب‌پذیری اسکن می‌کنند.

نمونه‌ها: \lr{scanssh, tiger, zmap}

\subsubsection{\lr{blackarch-sniffer}}
بسته‌هایی که به تحلیل ترافیک شبکه مربوط می‌شوند.

نمونه‌ها: \lr{hexinject, pytactle, xspy}

\subsubsection{\lr{blackarch-social}}
بسته‌هایی که عمدتاً به سایت‌های شبکه‌ی اجتماعی حمله می‌کنند.

نمونه‌ها: \lr{jigsaw, websploit}

\subsubsection{\lr{blackarch-spoof}}
بسته‌هایی که می‌کوشند هویت مهاجم را جعل کنند، به‌طوری که مهاجم در نظر قربانی
مهاجم دیده نشود.

نمونه‌ها: \lr{arpoison, lans, netcommander}

\subsubsection{\lr{blackarch-threat-model}}
بسته‌هایی که برای گزارش‌دهی یا ثبت مدل تهدید تعریف‌شده در یک سناریوی مشخص به کار
می‌روند.

نمونه‌ها: \lr{magictree}

\subsubsection{\lr{blackarch-tunnel}}
بسته‌هایی که برای تونل‌کردن ترافیک شبکه روی یک شبکه‌ی مشخص به کار می‌روند.

نمونه‌ها: \lr{ctunnel, iodine, ptunnel}

\subsubsection{\lr{blackarch-unpacker}}
بسته‌هایی که برای بیرون‌کشیدن بدافزارِ از پیش بسته‌بندی‌شده از یک فایل اجرایی
به کار می‌روند.

نمونه‌ها: \lr{js-beautify}

\subsubsection{\lr{blackarch-voip}}
بسته‌هایی که روی برنامه‌ها و پروتکل‌های \lr{VoIP} کار می‌کنند.

نمونه‌ها: \lr{iaxflood, rtp-flood, teardown}

\subsubsection{\lr{blackarch-webapp}}
بسته‌هایی که روی برنامه‌های در معرض اینترنت (\lr{internet-facing}) کار می‌کنند.

نمونه‌ها: \lr{metoscan, whatweb, zaproxy}

\subsubsection{\lr{blackarch-windows}}
این گروه برای هر بسته‌ی بومیِ ویندوز است که از طریق \lr{wine} اجرا می‌شود.

نمونه‌ها: \lr{3proxy-win32, pwdump, winexe}

\subsubsection{\lr{blackarch-wireless}}
بسته‌هایی که در هر سطحی روی شبکه‌های بی‌سیم کار می‌کنند.

نمونه‌ها: \lr{airpwn, mdk3, wiffy}

\section{ساختار مخزن}
مخزن گیت اصلی \lr{BlackArch} را اینجا می‌توانید پیدا کنید:
\href{https://github.com/BlackArch/blackarch}{\lr{https://github.com/BlackArch/blackarch}}.
چند مخزن فرعی هم اینجا هستند:
\href{https://github.com/BlackArch}{\lr{https://github.com/BlackArch}}.

در مخزن گیت اصلی سه پوشه‌ی مهم وجود دارد:

\begin{itemize}
\item \lr{docs}: مستندات.
\item \lr{packages}: فایل‌های \lr{PKGBUILD}.
\item \lr{scripts}: اسکریپت‌های کوچک و کاربردی.
\end{itemize}

\subsection{اسکریپت‌ها}
در اینجا مرجعی برای اسکریپت‌های پوشه‌ی \lr{\texttt{scripts}} آمده است:

\begin{itemize}
\item \lr{baaur}: به‌زودی بسته‌ها را در \lr{AUR} بارگذاری خواهد کرد.
\item \lr{babuild}: ساخت یک بسته.
\item \lr{bachroot}: مدیریت یک \lr{chroot} برای آزمایش.
\item \lr{baclean}: پاک‌کردن فایل‌های قدیمی بسته (با پسوند \lr{pkg.tar.xz}) از مخزن بسته‌ها.
\item \lr{baconflict}: به‌زودی جایگزین \lr{\texttt{conflicts}} خواهد شد.
\item \lr{bad-files}: یافتن فایل‌های نامطلوب در بسته‌های ساخته‌شده.
\item \lr{balock}: گرفتن یا آزادکردن قفل مخزن بسته‌ها.
\item \lr{banotify}: اطلاع‌دادن به \lr{IRC} درباره‌ی \lr{push} شدن بسته‌ها.
\item \lr{barelease}: انتشار بسته‌ها در مخزن بسته‌ها.
\item \lr{baright}: چاپ اطلاعات کپی‌رایت \lr{BlackArch}.
\item \lr{basign}: امضای بسته‌ها.
\item \lr{basign-key}: امضای یک کلید.
\item \lr{blackman}: تقریباً مثل \lr{pacman} رفتار می‌کند اما از روی \lr{git} می‌سازد
    (با \lr{Blackman} ساخته‌ی \lr{nrz} اشتباه گرفته نشود).
\item \lr{check-groups}: بررسی گروه‌ها.
\item \lr{checkpkgs}: بررسی بسته‌ها برای یافتن خطا.
\item \lr{conflicts}: بررسی تداخل فایل‌ها.
\item \lr{dbmod}: تغییر دادن پایگاه داده‌ی بسته‌ها.
\item \lr{depth-list}: ساخت فهرستی مرتب‌شده بر اساس عمق وابستگی.
\item \lr{deptree}: ساخت درخت وابستگی، با فهرست‌کردن فقط بسته‌هایی که \lr{BlackArch}
    ارائه می‌کند.
\item \lr{get-blackarch-deps}: گرفتن فهرست وابستگی‌های \lr{blackarch} یک بسته.
\item \lr{get-official}: گرفتن بسته‌های رسمی برای انتشار.
\item \lr{list-loose-packages}: فهرست‌کردن بسته‌هایی که در هیچ گروهی نیستند و
    وابستگی بسته‌ی دیگری هم نیستند.
\item \lr{list-needed}: فهرست‌کردن وابستگی‌های گم‌شده.
\item \lr{list-removed}: فهرست‌کردن بسته‌هایی که در مخزن بسته‌ها هستند اما در \lr{git}
    نیستند.
\item \lr{list-tools}: فهرست‌کردن ابزارها.
\item \lr{outdated}: یافتن بسته‌هایی که در مخزن بسته‌ها نسبت به مخزن \lr{git} قدیمی‌اند.
\item \lr{pkgmod}: تغییر یک بسته‌ی ساخته‌شده.
\item \lr{pkgrel}: افزایش \lr{pkgrel} در یک بسته.
\item \lr{prep}: مرتب‌کردن سبک (استایل) فایل \lr{PKGBUILD} و یافتن خطاها.
\item \lr{sitesync}: همگام‌سازی بین نسخه‌ی محلی مخزن بسته‌ها و نسخه‌ی راه‌دور.
\item \lr{size-hunt}: یافتن بسته‌های حجیم.
\item \lr{source-backup}: پشتیبان‌گیری از فایل‌های سورس بسته‌ها.
\end{itemize}

\section{مشارکت در مخزن}
این بخش نشان می‌دهد چگونه در پروژه‌ی \lr{BlackArch Linux} مشارکت کنید. ما
درخواست ادغام (\lr{pull request}) در هر اندازه‌ای را می‌پذیریم، از اصلاح یک غلط
تایپی کوچک تا افزودن بسته‌ی جدید.\\برای راهنمایی، پیشنهاد یا پرسش، آزادانه با
ما تماس بگیرید.
\\\\
از مشارکت همه استقبال می‌کنیم و قدردان همه‌ی مشارکت‌ها هستیم.

\subsection{آموزش‌های ضروری}
لطفاً پیش از مشارکت، آموزش‌های زیر را بخوانید:
\begin{itemize}
\item
\href{https://wiki.archlinux.org/index.php/Arch\_Packaging\_Standards)}{\lr{Arch Packaging Standards}}
\item \href{https://wiki.archlinux.org/index.php/Creating\_Packages}{\lr{Creating Packages}}
\item \href{https://wiki.archlinux.org/index.php/PKGBUILD}{\lr{PKGBUILD}}
\item \href{https://wiki.archlinux.org/index.php/Makepkg}{\lr{Makepkg}}
\end{itemize}

\subsection{مراحل مشارکت}
برای ارسال تغییراتتان به پروژه‌ی \lr{BlackArchLinux} این مراحل را دنبال کنید:
\begin{enumerate}
\item مخزن را از
\url{https://github.com/BlackArch/blackarch}
فورک (\lr{fork}) کنید.
\item فایل‌های لازم را ویرایش کنید (مثلاً \lr{PKGBUILD}، فایل‌های \lr{patch} و
غیره).
\item تغییراتتان را کامیت (\lr{commit}) کنید.
\item تغییراتتان را پوش (\lr{push}) کنید.
\item از ما بخواهید تغییراتتان را ادغام (\lr{merge}) کنیم؛ ترجیحاً از طریق
\lr{pull request}.
\end{enumerate}

\subsection{نمونه}
نمونه‌ی زیر نشان می‌دهد چگونه یک بسته‌ی جدید را برای پروژه‌ی \lr{BlackArch} ارسال
کنید. در اینجا با \href{https://github.com/Jguer/yay}{\lr{yay}}
(می‌توانید از \lr{pacaur} هم استفاده کنید) فایل \lr{PKGBUILD} موجود برای
\textbf{\lr{nfsshell}} را از \href{https://aur.archlinux.org/}{\lr{AUR}} می‌گیریم و
آن را مطابق نیازمان اصلاح می‌کنیم.

\subsubsection{دریافت \lr{PKGBUILD}}
فایل \textit{\lr{PKGBUILD}} را با \lr{yay} یا \lr{pacaur} دریافت کنید:
\begin{latin}
\begin{lstlisting}
  user@blackarchlinux $ yay -G nfsshell
  ==> Download nfsshell sources
  x LICENSE
  x PKGBUILD
  x gcc.patch
  user@blackarchlinux $ cd nfsshell/
\end{lstlisting}
\end{latin}

\subsubsection{مرتب‌کردن \lr{PKGBUILD}}
فایل \textit{\lr{PKGBUILD}} را مرتب کنید تا در وقتتان صرفه‌جویی شود:
\begin{latin}
\begin{lstlisting}
  user@blackarchlinux nfsshell $ ./blackarch/scripts/prep PKGBUILD
  cleaning 'PKGBUILD'...
  expanding tabs...
  removing vim modeline...
  removing id comment...
  removing contributor and maintainer comments...
  squeezing extra blank lines...
  removing '|| return'...
  removing leading blank line...
  removing $pkgname...
  removing trailing whitespace...
\end{lstlisting}
\end{latin}

\subsubsection{اصلاح \lr{PKGBUILD}}
فایل \textit{\lr{PKGBUILD}} را اصلاح کنید:
\begin{latin}
\begin{lstlisting}
  user@blackarchlinux nfsshell $ vi PKGBUILD
\end{lstlisting}
\end{latin}

\subsubsection{ساخت بسته}
بسته را بسازید:
\begin{latin}
\begin{lstlisting}
user@blackarchlinux nfsshell $ makepkg -sf
==> Making package: nfsshell 19980519-1 (Mon Dec  2 17:23:51 CET 2013)
==> Checking runtime dependencies...
==> Checking buildtime dependencies...
==> Retrieving sources...
-> Downloading nfsshell.tar.gz...
% Total    % Received % Xferd  Average Speed   Time    Time     Time
CurrentDload  Upload   Total   Spent    Left  Speed100 29213  100 29213    0
0  48150      0 --:--:-- --:--:-- --:--:-- 48206
-> Found gcc.patch
-> Found LICENSE
...
<lots of build process and compiler output here>
...
==> Leaving fakeroot environment.
==> Finished making: nfsshell 19980519-1 (Mon Dec  2 17:23:53 CET 2013)
\end{lstlisting}
\end{latin}

\subsubsection{نصب و آزمایش بسته}
بسته را نصب و آزمایش کنید:
\begin{latin}
\begin{lstlisting}
  user@blackarchlinux nfsshell $ pacman -U nfsshell-19980519-1-x86_64.pkg.tar.xz
  user@blackarchlinux nfsshell $ nfsshell # test it
\end{lstlisting}
\end{latin}

\subsubsection{افزودن، کامیت و پوش‌کردن بسته}
بسته را \lr{add} و \lr{commit} و \lr{push} کنید:
\begin{latin}
\begin{lstlisting}
user@blackarchlinux nfsshell $ cd /blackarchlinux/packages
user@blackarchlinux ~/blackarchlinux/packages $ mv ~/nfsshell .
user@blackarchlinux ~/blackarchlinux/packages $ git commit -am nfsshell && git push
\end{lstlisting}
\end{latin}

\subsubsection{ایجاد \lr{pull request}}
در \href{https://github.com/}{\lr{github.com}} یک \lr{pull request} بسازید:
\begin{latin}
\begin{lstlisting}
  firefox https://github.com/<contributor>/blackarchlinux
\end{lstlisting}
\end{latin}

\subsubsection{افزودن \lr{remote} برای \lr{upstream}}
اگر روی یک فورک کار می‌کنید و می‌خواهید با مخزن اصلی (\lr{upstream}) هماهنگ بمانید،
کار هوشمندانه این است که فورک خودتان را \lr{pull} کنید و مخزن اصلی \lr{BlackArch} را
به‌عنوان یک \lr{remote} اضافه کنید:
\begin{latin}
\begin{lstlisting}
  user@blackarchlinux ~/blackarchlinux $ git remote -v
  origin <the url of your fork> (fetch)
  origin <the url of your fork> (push)
  user@blackarchlinux ~/blackarchlinux $ git remote add upstream https://github.com/blackarch/blackarch
  user@blackarchlinux ~/blackarchlinux $ git remote -v
  origin <the url of your fork> (fetch)
  origin <the url of your fork> (push)
  upstream https://github.com/blackarch/blackarch (fetch)
  upstream https://github.com/blackarch/blackarch (push)
\end{lstlisting}
\end{latin}

به‌طور پیش‌فرض \lr{git} باید مستقیماً به \lr{origin} پوش کند، ولی مطمئن شوید تنظیمات \lr{git}
شما درست است. اگر دسترسی کامیت ندارید این موضوع مشکلی ایجاد نمی‌کند، چون بدون
آن دسترسی نمی‌توانید به \lr{upstream} پوش کنید.

اگر دسترسی کامیت دارید، شاید با
\lr{\textit{git@github.com:blackarch/blackarch.git}} نتیجه‌ی بهتری بگیرید،
ولی انتخاب با خودتان است.

\subsection{خواسته‌های ما}
\begin{enumerate}
\item در فایل‌های \textit{\lr{PKGBUILD}} کامنت \textbf{\lr{Maintainer}} یا
\textbf{\lr{Contributor}} اضافه نکنید. نام نگه‌دارندگان (\lr{maintainer}) و
مشارکت‌کنندگان را به بخش «مشارکت‌کنندگان» (\lr{AUTHORS}) راهنمای \lr{BlackArch} اضافه
کنید.
\item برای یکدست‌بودن، لطفاً از سبک کلی سایر فایل‌های \textit{\lr{PKGBUILD}} مخزن
پیروی کنید و تورفتگی را دو فاصله (\lr{space}) بگیرید.
\end{enumerate}

\subsection{نکات عمومی}
ابزار \href{http://wiki.archlinux.org/index.php/Namcap}{\lr{namcap}} می‌تواند
بسته‌ها را از نظر خطا بررسی کند.

%------------------%
%  Chapter 4       %
%------------------%

\chapter{راهنمای ابزارها}
به‌زودی...

\section{به‌زودی}
به‌زودی...

%%% APPENDIX %%%
\appendix
%%%%%%%%%%%%%%%%%%%%%%%%%%%%%%%%%%%%%%%%%%%%%%%%%%%%%%%%%%%%%%%%%%%%%%%%%%%%%%%%
%                                                                              %
% BlackArch Linux appendix template (Persian / fa)                             %
%                                                                              %
%%%%%%%%%%%%%%%%%%%%%%%%%%%%%%%%%%%%%%%%%%%%%%%%%%%%%%%%%%%%%%%%%%%%%%%%%%%%%%%%

\appendix

\chapter{پیوست}

\section{پرسش‌های متداول}

\section{مشارکت‌کنندگان}
افراد زیر مستقیماً در \lr{BlackArch} مشارکت داشته‌اند:
% names and e-mail addresses are Latin text: typeset them left-to-right with the Latin font
\begin{latin}
\begin{itemize}
\item Tyler Bennnett (tylerb@trix2voip.com)
\item fnord0 (fnord0@riseup.net)
\item nrz (nrz@nullsecurity.net)
\item Ellis Kenyo (elken.tdos@gmail.com)
\item CaledoniaProject (the.warl0ck.1989@gmail.com)
\item sudokode (sudokode@gmail.com)
\item Valentin Churavy (v.churavy@gmail.com)
\item Boy Sandy Gladies Arriezona (reno.esper@gmail.com)
\item Mathias Nyman
\item Johannes Löthberg (demizide@gmail.com)
\item Thiago da Silva Teixeira (teixeira.zeus@gmail.com)
\end{itemize}
\end{latin}
افراد زیر مستقیماً در \lr{ArchPwn} مشارکت داشته‌اند
و سپس به \lr{BlackArch} پیوسته‌اند:
\begin{latin}
\begin{itemize}
\item Francesco Piccinno (stack.box@gmail.com)
\item jensp (jens@jenux.homelinux.org)
\item Valentin Churavy (v.churavy@gmail.com)
\end{itemize}
\end{latin}
برای ساخت بسته‌ها از کدهای افراد زیر استفاده کرده‌ایم:
\begin{latin}
\begin{itemize}
\item  3ED (krzysztof1987@gmail.com)
\item  AUR Perl (aurperl@juster.info)
\item  Aaron Griffin (aaron@archlinux.org)
\item  Abakus (java5@arcor.de)
\item  Adam Wolk (netprobe@gmail.com)
\item  Aleix Pol (aleixpol@kde.org)
\item  Aleshus (aleshusi@gmail.com)
\item  Alessandro Pazzaglia (jackdroido@gmail.com)
\item  Alessandro Sagratini (ale\_sagra@hotmail.com)
\item  Alex Cartwright (alexc223@googlemail.com)
\item  Alexander De Sousa (archaur.xandy21@spamgourmet.com)
\item  Alexander Rødseth (rodseth@gmail.com)
\item  Allan McRae (allan@archlinux.org)
\item  AmaN (gabroo.punjab.da@gmail.com)
\item  Andre Klitzing (aklitzing@online.de)
\item  Andrea Scarpino (andrea@archlinux.org)
\item  Andreas Schönfelder (passtschu@freenet.de)
\item  Andrej Gelenberg (andrej.gelenberg@udo.edu)
\item  Angel Velasquez (angvp@archlinux.org)
\item  Antoine Lubineau (antoine@lubignon.info)
\item  Anton Bazhenov (anton.bazhenov@gmail.com)
\item  Arkham (arkham@archlinux.us)
\item  Arthur Danskin (arthurdanskin@gmail.com)
\item  Balda (balda@balda.ch)
\item  Balló György (ballogyor+arch@gmail.com)
\item  Bartek Piotrowski (barthalion@gmail.com)
\item  Bartosz Feński (fenio@debian.org)
\item  Bartłomiej Piotrowski (nospam@bpiotrowski.pl)
\item  Bogdan Szczurek (thebodzio@gmail.com)
\item  Brad Fanella (bradfanella@archlinux.us)
\item  Brian Bidulock (bidulock@openss7.org)
\item  C Anthony Risinger (anthony@xtfx.me)
\item  CRT (crt.011@gmail.com)
\item  Can Celasun (dcelasun@gmail.com)
\item  Chaniyth (chaniyth@yahoo.com)
\item  Chris Brannon (cmbrannon79@gmail.com)
\item  Chris Giles (Chris.G.27@gmail.com) \& daschu117
\item  Christoph Siegenthaler (csi@gmx.ch)
\item  Christoph Zeiler (archNOSPAM@moonblade.org)
\item  Clément DEMOULINS (clement@archivel.fr)
\item  Corrado Primier (bardo@aur.archlinux.org)
\item  Daenyth (Daenyth+Arch@gmail.com)
\item  Dale Blount (dale@archlinux.org)
\item  Damir Perisa (damir.perisa@bluewin.ch)
\item  Dan Fuhry (dan@fuhry.us)
\item  Dan Serban (dserban01@yahoo.com)
\item  Daniel A. Campoverde Carrión
\item  Daniel Golle
\item  Daniel Griffiths (ghost1227@archlinux.us)
\item  Daniel J Griffiths (ghost1227@archlinux.us)
\item  Daniel Micay (danielmicay@gmail.com)
\item  Dave Reisner (dreisner@archlinux.org)
\item  Dawid Wrobel (cromo@klej.net)
\item  Devaev Maxim (mdevaev@gmail.com)
\item  Devin Cofer (ranguvar@archlinux.us)
\item  DigitalPathogen (aur@InfoSecResearchLabs.co.uk)
\item  DigitalPathogen (aur@digitalpathogen.co.uk)
\item  Dmitry A. Ilyashevich (dmitry.ilyashevich@gmail.com)
\item  Dominik Heidler (dheidler@gmail.com)
\item  DrZaius (lou@fakeoutdoorsman.com)
\item  Ebubekir KARUL (ebubekirkarul@yandex.com)
\item  Eduard "bekks" Warkentin (eduard.warkentin@gmail.com)
\item  Elmo Todurov (todurov@gmail.com)
\item  Emmanuel Gil Peyrot (linkmauve@linkmauve.fr)
\item  Eric Belanger (eric@archlinux.org)
\item  Ermak (ermak@email.it)
\item  Evangelos Foutras (evangelos@foutrelis.com)
\item  Fabian Melters (melters@gmail.com)
\item  Fabiano Furtado (fusca14@gmail.com)
\item  Federico Quagliata (ntp@quaqo.org)
\item  Firmicus (francois.archlinux@org)
\item  Florian Pritz (bluewind@jabber.ccc.de)
\item  Florian Pritz (flo@xinu.at)
\item  Francesco Piccinno (stack.box@gmail.com)
\item  François Charette (francois@archlinux.org)
\item  Gaetan Bisson (bisson@archlinux.org)
\item  Geoffroy Carrier (geoffroy.carrier@koon.fr)
\item  Georg Grabler (STiAT)
\item  George Hilliard (gh403@msstate.edu)
\item  Gerardo Exequiel Pozzi (vmlinuz386@yahoo.com.ar)
\item  Gilles CHAUVIN (gcnweb@gmail.com)
\item  Giovanni Scafora (giovanni@archlinux.org)
\item  Gordin (9ordin@gmail.com)
\item  Guillaume ALAUX (guillaume@archlinux.org)
\item  Guillermo Ramos (0xwille@gmail.com)
\item  Gustavo Alvarez (sl1pkn07@gmail.com)
\item  Hugo Doria (hugo@archlinux.org)
\item  Hyacinthe Cartiaux (hyacinthe.cartiaux@free.fr)
\item  James Fryman (jfryman@gmail.com)
\item  Jan "heftig" Steffens (jan.steffens@gmail.com)
\item  Jan de Groot (jgc@archlinux.org)
\item  Jaroslav Lichtblau (dragonlord@aur.archlinux.org)
\item  Jaroslaw Swierczynski (swiergot@aur.archlinux.org)
\item  Jason Chu (jason@archlinux.org)
\item  Jason R Begley (jayray@digitalgoat.com)
\item  Jason Rodriguez
\item  Jason St. John (jstjohn@purdue.edu)
\item  Jawmare (victor2008@gmail.com)
\item  Jeff Mickey (jeff@archlinux.org)
\item  Jens Pranaitis (jens@chaox.net)
\item  Jens Pranaitis (jens@jenux.homelinux.org)
\item  Jinx (jinxware@gmail.com)
\item  John D Jones III (jnbek1972@gmail.com)
\item  John Proctor (jproctor@prium.net)
\item  Jon Bergli Heier (snakebite@jvnv.net)
\item  Jonas Heinrich
\item  Jonathan Steel (jsteel@aur.archlinux.org)
\item  Joris Steyn (jorissteyn@gmail.com)
\item  Josh VanderLinden (arch@cloudlery.com)
\item  Jozef Riha (jose1711@gmail.com)
\item  Judd Vinet (jvinet@zeroflux.org)
\item  Juergen Hoetzel (jason@archlinux.org)
\item  Juergen Hoetzel (juergen@archlinux.org)
\item  Justin Davis (jrcd83@gmail.com)
\item  Kaiting Chen (kaitocracy@gmail.com)
\item  Kaos
\item  Kevin Piche (kevin@archlinux.org)
\item  Kory Woods (kory@virlo.net)
\item  Kyle Keen (keenerd@gmail.com)
\item  Larry Hajali (larryhaja@gmail.com)
\item  LeCrayonVert
\item  Le\_suisse (lesuisse.dev+aur@gmail.com)
\item  Lekensteyn (lekensteyn@gmail.com)
\item  Limao Luo (luolimao+AUR@gmail.com)
\item  Lucien Immink
\item  Lukas Fleischer (archlinux@cryptocrack.de)
\item  Manolis Tzanidakis
\item  Marcin "avalan" Falkiewicz (avalatron@gmail.com)
\item  Mariano Verdu (verdumariano@gmail.com)
\item  Marti Raudsepp (marti@juffo.org)
\item  MatToufoutu (mattoufootu@gmail.com)
\item  Matthew Sharpe (matt.sharpe@gmail.com)
\item  Mauro Andreolini (mauro.andreolini@unimore.it)
\item  Max Pray a.k.a. Synthead (synthead@gmail.com)
\item  Max Roder (maxroder@web.de)
\item  Maxwell Pray a.k.a. Synthead (synthead@gmail.com)
\item  Maxwell Pray a.k.a. Synthead (synthead1@gmail.com)
\item  Mech (tiago.bmp@gmail.com)
\item  Michael Düll (mail@akurei.me)
\item  Michael P (ptchinster@archlinux.us)
\item  Michal Krenek (mikos@sg1.cz)
\item  Michal Zalewski (lcamtuf@coredump.cx)
\item  Miguel Paolino (mpaolino@gmail.com)
\item  Miguel Revilla (yo@miguelrevilla.com)
\item  Mike Roberts (noodlesgc@gmail.com)
\item  Mike Sampson (mike@sambodata.com)
\item  Nassim Kacha (nassim.kacha@gmail.com)
\item  Nicolas Pouillard (nicolas.pouillard@gmail.com)
\item  Nicolas Pouillard https://nicolaspouillard.fr
\item  Niklas Schmuecker
\item  Oleander Reis (oleander@oleander.cc)
\item  Olivier Le Moal (mail@olivierlemoal.fr)
\item  Olivier Médoc "oliv" (o\_medoc@yahoo.fr)
\item  Pascal E. (archlinux@hardfalcon.net)
\item  Patrick Leslie Polzer (leslie.polzer@gmx.net)
\item  Paul Mattal (paul@archlinux.org)
\item  Paul Mattal (pjmattal@elys.com)
\item  Pengyu CHEN (cpy.prefers.you@gmail.com)
\item  Peter Wu (lekensteyn@gmail.com)
\item  Philipp 'TamCore' B. (philipp@tamcore.eu)
\item  Pierre Schmitz (pierre@archlinux.de)
\item  Pranay Kanwar (pranay.kanwar@gmail.com)
\item  Pranay Kanwar (warl0ck@metaeye.org)
\item  PyroPeter (googlemail@com.abi1789)
\item  PyroPeter (googlemail.com@abi1789)
\item  Ray Rashif (schiv@archlinux.org)
\item  Remi Gacogne
\item  Renan Fernandes (renan@kauamanga.com)
\item  Richard Murri (admin@richardmurri.com)
\item  Roberto Alsina (ralsina@kde.org)
\item  Robson Peixoto (robsonpeixoto@gmail.com)
\item  Roel Blaauwgeers (roel@ttys0.nl)
\item  Rorschach (r0rschach@lavabit.com)
\item  Ruben Schuller (shiml@orgizm.net)
\item  Rudy Matela (rudy@matela.com)
\item  Ryon Sherman (ryon.sherman@gmail.com)
\item  Sabart Otto \item  Seberm (seberm@gmail.com)
\item  SakalisC (chrissakalis@gmail.com)
\item  Sam Stuewe (halosghost@archlinux.info)
\item  SanskritFritz (SanskritFritz@gmail.com)
\item  Sarah Hay (sarahhay@mb.sympatico)
\item  Sebastian Benvenuti (sebastianbenvenuti@gmail.com)
\item  Sebastian Nowicki (sebnow@gmail.com)
\item  Sebastien Duquette (ekse.0x@gmail.com)
\item  Sebastien LEDUC (sebastien@sleduc.fr)
\item  Sebastien Leduc (sebastien@sleduc.fr)
\item  Sergej Pupykin (pupykin.s+arch@gmail.com)
\item  Sergio Rubio (rubiojr@biondofu.net)
\item  Sheng Yu (magicfish1990@gmail.com)
\item  Simon Busch (morphis@gravedo.de)
\item  Simon Legner (Simon.Legner@gmail.com)
\item  Sirat18 (aur@sirat18.de)
\item  SpepS (dreamspepser@yahoo.it)
\item  Spider.007 (archPackage@spider007.net)
\item  Stefan Seering
\item  Stephane Travostino (stephane.travostino@gmail.com)
\item  Stéphane Gaudreault (stephane@archlinux.org)
\item  Sven Kauber (celeon@gmail.com)
\item  Sven Schulz (omee@archlinux.de)
\item  Sébastien Duquette (ekse.0x@gmail.com)
\item  Sébastien Luttringer (seblu@archlinux.org)
\item  TDY (tdy@gmx.com)
\item  Teemu Rytilahti (tpr@iki.fi)
\item  Testuser\_01
\item  Thanx (thanxm@gmail.com)
\item  Thayer Williams (thayer@archlinux.org)
\item  Thomas S Hatch (thatch45@gmail.com)
\item  Thorsten Töpper
\item  Tilmann Becker (tilmann.becker@freenet.de)
\item  Timothy Redaelli (timothy.redaelli@gmail.com)
\item  Timothée Ravier (tim@siosm.fr)
\item  Tino Reichardt
\item  Tobias Kieslich (tobias@justdreams.de)
\item  Tobias Powalowski (tpowa@archlinux.org)
\item  Tom K (tomk@runbox.com)
\item  Tom Newsom (Jeepster@gmx.co.uk)
\item  Tomas Lindquist Olsen (tomas@famolsen.dk)
\item  Travis Willard (travisw@wmpub.ca)
\item  Valentin Churavy (v.churavy@gmail.com)
\item  ViNS (gladiator@fastwebnet.it)
\item  Vlatko Kosturjak (kost@linux.hr)
\item  Wes Brown (wesbrown18@gmail.com)
\item  William Rea (sillywilly@gmail.com)
\item  Xavier Devlamynck (magicrhesus@ouranos.be)
\item  Xiao\item Long Chen (chenxiaolong@cxl.epac.to)
\item  aeolist (aeolist@hotmail.com)
\item  ality@pbrane.org
\item  astaroth (astaroth\_@web.de)
\item  bender02@archlinux.us
\item  billycongo (billycongo@gmail.com)
\item  bslackr (brendan@vastactive.com)
\item  cbreaker (cbreaker@tlen.pl)
\item  chimeracoder (dev@chimeracoder.net)
\item  damir (damir@archlinux.org)
\item  danitool
\item  darkapex (me@jailuthra.in)
\item  daronin
\item  dkaylor (dpkaylor@gmail.com)
\item  dobo (dobo90\_at\_gmail@com)
\item  dorphell (dorphell@archlinux.org)
\item  evr (evanroman@at.gmail)
\item  fnord0 (fnord0@riseup.net)
\item  fxbru (frxbru@gmail)
\item  hcar
\item  icarus (icarus.roaming@gmail.com)
\item  iceman (icemanf@gmail.com)
\item  kastor (kastor@fobos.org)
\item  kfgz (kfgz@interia.pl)
\item  linuxSEAT (linuxSEAT@gmail.com)
\item  m4xm4n (max@maxfierke.com)
\item  mar77i (mysatyre@gmail.com)
\item  marc0s (marc0s@fsfe.org)
\item  mickael9 (mickael9@gmail.com)
\item  nblock (nblock@archlinux.us)
\item  nofxx (x@nofxx.com)
\item  onny (onny@project
\item  pootzko (pootzko@gmail.com)
\item  revel (revel@muub.net)
\item  rich\_o (rich\_o@lavabit.com)
\item  s1gma (s1gma@mindslicer.com)
\item  sandman (r.coded@gmail.com)
\item  sebikul (sebikul@gmail.com)
\item  sh0 (mee@sh0.org)
\item  shild (sxp@bk.ru)
\item  simo (simo@archlinux.org)
\item  snuo
\item  sudokode (sudokode@gmail.com)
\item  tobias (tobias@archlinux.org)
\item  trashstar (trash@ps3zone.org)
\item  unexist (unexist@subforge.org)
\item  untitled (rnd0x00@gmail.com)
\item  virtuemood (virtue@archlinux.us)
\item  wido (widomaker2k7@gmail.com)
\item  wodim (neikokz@gmail.com)
\item  yannsen (ynnsen@gmail.com)
\end{itemize}
\end{latin}


\end{document}

%%% EOF %%%
